\section{Related Work}
\label{sec:related}

Legal NLP has studied judgment prediction, statutory reasoning, legal benchmarks, case retrieval and entailment, and summarization, often with domain-adapted encoders.
Audits show that LLMs and commercial legal RAG tools still hallucinate.
Prior work also studies overruling detection as a classification task over case text; we instead extract treatment edges to build time-indexed validity intervals and apply them as a retrieval gate.
These lines of work take legal text as the premise, treat the corpus as static, and do not test whether retrieved authority remained operative at a given date.
\plre{} changes the inference unit: product attributes are mapped to legal-risk theories in a directional, product-conditioned, and temporal manner.

Time-sensitive QA conditions answers on the moment a question is asked; temporal IR uses document time as a ranking signal.
Both treat time as a property of facts or queries.
In \plre{}, time enters through supersession: a later document alters the evidential status of an earlier one.
Commercial citators track negative treatment; citation-network analysis and edge labeling characterize the underlying graph.
\sysname{} attaches validity intervals to extracted treatment edges and applies them as a retrieval-time gate.

RAG typically executes a fixed retrieval plan before generation; work on knowledge conflicts and outdated parametric knowledge shows that generators may prefer stale but familiar evidence unless retrieval is constrained.
Agent frameworks benefit from iterative tool use yet stop when the model itself deems its answer satisfactory.
\sysname{} ties stopping to an externally inspectable predicate over evidence support and mitigation coverage, and employs conformal risk control narrowly for evidence omission rather than legal correctness.
